\section{Probability laws and identifying transforms}
\label{final:probability}

Throughout this section we use the intrinsic decomposition of $\Sigma$
in its marked positive coordinate:
\[
 H(t)=\log t-t+1,\qquad I(t)=t-1-\log t,\qquad
 p(t)=t e^{-t}.
\]
The symbol $\overline F$ denotes a survival function.
All Gamma parameters are shape and
\emph{rate}, in that order.  Equality of probability laws identifies
densities almost everywhere; a pointwise density assertion uses their
specified continuous representatives.  Placement parameters do not
occur in these intrinsic statements.

\begin{proposition}[Uniqueness principles for the transforms used below]
\label{final:P0-uniqueness}
Finite Borel measures on $\mathbb R$ are determined by their
characteristic functions.  Finite measures on $[0,\infty)$ are
determined by their Laplace transforms on $[0,\infty)$, and finite
measures on $[0,1]$ by all their nonnegative integer moments.
Two probability laws with moment generating functions finite on
an open interval about zero and equal there are equal.
\end{proposition}
\begin{proof}
For the compact-interval assertion, equality on monomials gives
equality on polynomials, then on continuous functions by uniform
polynomial approximation, and finally on Borel sets by uniqueness
of finite regular measures.  Pushing a measure on $[0,\infty)$ forward
by $x\mapsto e^{-x}$ reduces Laplace uniqueness to this assertion:
integer Laplace values give all moments of the pushforward on $[0,1]$.

For characteristic functions, subtract the two finite measures.
Convolution with a centered Gaussian of variance $\varepsilon>0$
has density whose Fourier transform is the vanishing characteristic
function times $e^{-\varepsilon\xi^2/2}$.  Fourier inversion, justified
by the integrable Gaussian factor and Fubini's theorem, makes this
convolution zero.  Gaussian approximate identities converge against
every continuous function of compact support, so the original signed
measure is zero.  Finally an exponential moment gives a holomorphic
bilateral transform on a vertical strip about the imaginary axis.
The identity theorem extends equality near zero to that strip;
characteristic-function uniqueness applies on the imaginary axis.
\end{proof}

\begin{theorem}[Gamma law and complete transform characterizations]
\label{final:P1}
Let $\mu_\Gamma(dt)=\mathbf 1_{(0,\infty)}(t)p(t)\,dt$.
It is the probability law $\operatorname{Gamma}(2,1)$, with topological
support $[0,\infty)$ and no atom at zero.  Its transforms are
\begin{align}
 \int e^{-zt}\,\mu_\Gamma(dt)&=(1+z)^{-2},
       &&\operatorname{Re}z>-1,\label{final:eq-gamma-laplace}\\
 \int e^{i\xi t}\,\mu_\Gamma(dt)&=(1-i\xi)^{-2},
       &&\xi\in\mathbb R,\label{final:eq-gamma-characteristic}\\
 \int t^{z-1}\,\mu_\Gamma(dt)&=\Gamma(z+1),
       &&\operatorname{Re}z>-1,\label{final:eq-gamma-mellin}\\
 \int t^n\,\mu_\Gamma(dt)&=(n+1)!,
       &&n\in\mathbb N_0.\label{final:eq-gamma-moments}
\end{align}
Each of the following complete observations identifies its candidate
as $\mu_\Gamma$:
\begin{enumerate}
\item a finite positive Borel measure on $\mathbb R$ with the moments
      \eqref{final:eq-gamma-moments}, including $n=0$;
\item a probability on $[0,\infty)$ with Laplace transform
      $(1+\lambda)^{-2}$ for every $\lambda\geq0$;
\item a probability on $\mathbb R$ with characteristic function
      $(1-i\xi)^{-2}$ for every real $\xi$;
\item a probability on $(0,\infty)$ with
      $\mathbb E T^{i\xi}=\Gamma(2+i\xi)$ for every real $\xi$.
\end{enumerate}
In the first and third clauses, positive support is a conclusion.
Moreover
\[
 \mathbb ET=2,\quad \operatorname{Var}(T)=2,\quad
 \mathbb E\log T=1-\gamma_E,\quad
 \kappa_n(T)=2(n-1)!\quad(n\geq1).
\]
\end{theorem}
\begin{proof}
Euler's integral and integration by parts give
$\int_0^\infty t^a e^{-bt}\,dt=\Gamma(a+1)b^{-a-1}$
for $a>-1$, $b>0$, with its holomorphic extension on the stated
half-planes.  Taking $a=1$ proves normalization and the first two
transforms; taking $a=z$ proves the Mellin formula and the moments.
Differentiating the Gamma integral at $a=1$ is justified by an
integrable logarithmic majorant near zero and exponential decay at
infinity.  It gives $\mathbb E\log T=\Gamma'(2)=1-\gamma_E$.
The identity $\log\mathbb E e^{uT}=-2\log(1-u)$ yields the cumulants.

For the first converse the zeroth moment gives mass one.  For
$0<c<1$, positivity and monotone convergence give
\[
 \int\cosh(ct)\,\mu(dt)
 =\sum_{n=0}^\infty(2n+1)c^{2n}
 =\frac{1+c^2}{(1-c^2)^2}<\infty.
\]
Since $e^{c|t|}\leq2\cosh(ct)$, the candidate has a two-sided
exponential moment.  Dominated power-series expansion for $|u|<1$
therefore gives
\[
 \int e^{ut}\,\mu(dt)
 =\sum_{n=0}^\infty\frac{u^n(n+1)!}{n!}
 =(1-u)^{-2}.
\]
Proposition~\ref{final:P0-uniqueness} identifies the measure on the
whole real line.  Its support follows from the strictly positive
density on $(0,\infty)$ and zero density on $(-\infty,0]$.
The second and third converses are the same uniqueness proposition.
In the fourth, the supplied function is the characteristic function
of $\log T$; identify that law and push it forward by the exponential.
\end{proof}

\begin{proposition}[Exponential samples also derive support]
\label{final:P1-samples}
If a finite positive Borel measure $\nu$ on $\mathbb R$ satisfies
\[
 \int e^{-ns}\,\nu(ds)=(n+1)^{-2}\qquad(n\in\mathbb N_0),
\]
then $\nu(ds)=\mathbf 1_{(0,\infty)}(s)s e^{-s}\,ds$.
\end{proposition}
\begin{proof}
The sample at zero gives mass one.  If
$\nu(( -\infty,-\varepsilon])=m>0$, then the $n$th sample is
at least $m e^{n\varepsilon}$, a contradiction as $n\to\infty$.
Taking a countable union excludes the entire negative ray.
Push forward by $y=e^{-s}$ to $[0,1]$.  Its moments are
$(n+1)^{-2}$, which are those of $(-\log y)\,dy$ on $(0,1)$.
Compact moment uniqueness and the inverse change of variables give
the stated law, including absence of an atom at zero.
\end{proof}

\begin{proposition}[Boundaries of moment identification]
\label{final:P1-boundaries}
Removing positivity, removing the zeroth moment, or retaining only
finitely many ordinary moments invalidates the corresponding
unrestricted moment converse in Theorem~\ref{final:P1}.
\end{proposition}
\begin{proof}
Without the zeroth equation, every $\mu_\Gamma+a\delta_0$, $a>0$,
has the specified moments of positive orders.  To remove positivity,
take a nonzero real even smooth Fourier bump supported in
$[-2,-1]\cup[1,2]$ and let $h$ be its inverse Fourier transform.
The real Schwartz function $h$ has every polynomial moment zero,
because those moments are constant multiples of derivatives of its
Fourier transform at zero.  Thus the distinct finite signed measures
$\mu_\Gamma+\varepsilon h(t)\,dt$ have all the same moments.

For any fixed $m$, choose $m+2$ nonzero smooth bumps with disjoint
compact supports in $(1,2)$.  The $m+1$ linear equations making a
linear combination $h$ orthogonal to $1,t,\ldots,t^m$ have a nonzero
solution.  Disjoint supports ensure $h\neq0$.  On this compact set
$p$ has a positive minimum, so $p+\varepsilon h$ is a positive
probability density for all sufficiently small $|\varepsilon|$.
It is distinct from $p$ and retains every requested moment.
\end{proof}

\begin{theorem}[Marked Poisson recurrence, cumulants and divergences]
\label{final:P2}
A probability sequence $(\pi_n)_{n\geq0}$ satisfies
\[
 (n+1)\pi_{n+1}=\pi_n\qquad(n\geq0)
\]
if and only if $\pi_n=e^{-1}/n!$.  For $N$ with this law,
\begin{align}
 K(u)&=\log\mathbb E e^{uN}=e^u-1,\label{final:eq-poisson-cgf}\\
 \Lambda(u)&=\log\mathbb E e^{u(N-1)}=e^u-1-u=I(e^u).
       \label{final:eq-poisson-centered}
\end{align}
The complete uncentered CGF near zero identifies $N$ as
$\operatorname{Pois}(1)$.  The complete centered expression identifies
the law of $N-1$ when the shift is marked as one.  As a scalar
function it reconstructs $I(t)=\Lambda(\log t)$.

With $D$ denoting relative entropy, for $a,b>0$ one has
\begin{align}
 D(\operatorname{Pois}(a)\Vert\operatorname{Pois}(b))
     &=a I(b/a),\label{final:eq-poisson-kl}\\
 D(\operatorname{Exp}(a)\Vert\operatorname{Exp}(b))
     &=I(b/a),\label{final:eq-exp-kl}\\
 2D(\mathcal N(0,a)\Vert\mathcal N(0,b))
     &=I(a/b),\label{final:eq-gaussian-kl}
\end{align}
where $a,b$ in the last line are variances.  Consequently each full
marked slice obtained by $(a,b)=(1,t),(1,t),(t,1)$, respectively,
reconstructs $I$.  With $\Psi(t)=t\log t-t+1$,
\[
 D(\operatorname{Pois}(a)\Vert\operatorname{Pois}(b))
 =\Psi(a)-\Psi(b)-\Psi'(b)(a-b).
\]
The Legendre transform of the centered CGF is
\[
 \Lambda^*(z)=
 \begin{cases}
 (1+z)\log(1+z)-z,&z>-1,\\
 1,&z=-1,\\
 +\infty,&z<-1.
 \end{cases}
\]
\end{theorem}
\begin{proof}
The recurrence gives $\pi_n=\pi_0/n!$ by induction; summing fixes
$\pi_0=e^{-1}$.  Conversely these coefficients satisfy the recurrence.
Summing the exponential series gives
$\mathbb E e^{uN}=e^{-1}\exp(e^u)$ and both CGFs.  Their converses
follow from Proposition~\ref{final:P0-uniqueness} applied to the actual
variable whose CGF is supplied.

For Poisson laws the log likelihood ratio at $n$ is
$b-a+n\log(a/b)$; its expectation is
$b-a+a\log(a/b)$.  For exponential laws the log ratio is
$\log(a/b)+(b-a)t$, whose expectation under rate $a$ is
$\log(a/b)+(b-a)/a$.  For centered Gaussians the log ratio is
$\frac12\log(b/a)+\frac12t^2(1/b-1/a)$, whose expectation uses
$\mathbb E t^2=a$.  Substitution proves the three formulas and the
Bregman identity.

For $z>-1$, maximizing $uz-\Lambda(u)$ gives $e^u=1+z$ and
the displayed value.  At $z=-1$ the objective is $1-e^u$ with
supremum one as $u\to-\infty$.  For $z<-1$ it tends to infinity
in that limit.  Exponential tilting by $uN$ gives
$\operatorname{Pois}(e^u)$, whose mean, variance, and Fisher
information in the $u$ coordinate are all $e^u$, as follows by
differentiating $K$ and by the score $N-e^u$.
\end{proof}

\begin{proposition}[The centering mark and KL orientation are essential]
\label{final:P2-boundaries}
An unmarked centered CGF does not determine the original location,
even for a law on $\mathbb N_0$.  Reversing the displayed divergence
slices changes their scalar potential.
\end{proposition}
\begin{proof}
Let $Y\sim\operatorname{Pois}(1)$ and $N_k=Y+k$ for an integer
$k\geq1$.  Then $N_k-\mathbb EN_k=Y-1$, with centered CGF
$e^u-1-u$, although $N_k$ is not $\operatorname{Pois}(1)$.
The reversed Poisson slice is $t\log t-t+1$; the reversed exponential
slice is $\log t+t^{-1}-1=I(1/t)$, and twice the reversed Gaussian
slice has this same value.  These follow directly from
\eqref{final:eq-poisson-kl}--\eqref{final:eq-gaussian-kl} and differ
from $I(t)$ in general.
\end{proof}

\begin{theorem}[Two marked max laws characterize the Gumbel CDF]
\label{final:P3}
Let $F$ be a CDF on $\mathbb R$.  The identities
\[
 F(x+\log2)^2=F(x),\qquad F(x+\log3)^3=F(x)
       \quad(x\in\mathbb R),\qquad F(0)=e^{-1}
\]
hold if and only if
\[
 F(x)=\exp(-e^{-x}).
\]
No density, smoothness, strict monotonicity, or full-support
hypothesis is required.  This CDF has density
$g(x)=e^{-x-e^{-x}}=p(e^{-x})$, and it satisfies the analogous
max identity for every positive integer $n$.
\end{theorem}
\begin{proof}
Iteration in both directions gives
$F(k\log2)=\exp(-2^{-k})$ for every integer $k$.
Monotone bounds by adjacent such values yield $0<F(x)<1$ at every
finite $x$.  Hence $G=-\log F$ is finite, positive, and nonincreasing.
The two equations imply
\[
 G(m\log2+n\log3)=e^{-(m\log2+n\log3)}
          \qquad(m,n\in\mathbb Z).
\]
The subgroup generated by $\log2$ and $\log3$ is dense in
$\mathbb R$: otherwise their ratio would be rational, contradicting
unique factorization in the equality $2^r=3^s$ for positive integers
$r,s$.  Approximating an arbitrary $x$ from both sides by subgroup
points and using monotonicity proves $G(x)=e^{-x}$.
Substitution proves the converse, the density formula, and the
max identities.  For independent observations the CDF of their
maximum is $F^n$, so these are precisely the stated shifted max laws.
\end{proof}

\begin{proposition}[Haar weighting and Gumbel calibration boundaries]
\label{final:P3-haar}
For the marked map $\tau(x)=e^{-x}$, the Gumbel probability satisfies
\[
 \tau_*(dF)=e^{-t}\,dt=p(t)\frac{dt}{t},\qquad
 \int_0^\infty p(t)\frac{dt}{t}=1.
\]
The measure $dt/t$ is multiplicative Haar measure and has infinite
mass.  Thus the Haar-weighted probability, and its complete density
in the marked logarithmic coordinate, each reconstruct $p$.
Removing the anchor in Theorem~\ref{final:P3}, or keeping only its
base-two equation, destroys uniqueness.
\end{proposition}
\begin{proof}
The positive change-of-variables Jacobian is $|dx/dt|=1/t$.
Since $g(-\log t)=p(t)$, it gives the pushforward formula;
conversely $p(t)=g(-\log t)$ recovers the density.
Invariance of $dt/t$ under $t\mapsto at$ is immediate by
substitution, and its mass diverges at both endpoints.

Without the anchor, $F_c(x)=\exp(-c e^{-x})$, $c>0$, satisfies
both max equations.  For only the base-two equation, put
$a=\log2$ and $h(x)=1+\varepsilon\sin(2\pi x/a)$.
For nonzero $\varepsilon$ sufficiently small, $h>0$ and $h'<h$.
Then $\exp(-e^{-x}h(x))$ is a strictly increasing CDF with the
correct endpoint limits and anchor; periodicity gives the base-two
equation.  It is not the Gumbel CDF.  The original Gamma probability
$p(t)\,dt$ and the Haar-weighted probability $p(t)\,dt/t$ are
different measures.
\end{proof}

\begin{proposition}[Canonical deficit law and coordinate involution]
\label{final:P4-data}
For $T\sim\mu_\Gamma$ let $V=I(T)$.  For $v\geq0$ let
$a(v)\leq1\leq b(v)$ be the roots of $I(t)=v$; explicitly
\[
 a(v)=-W_0(-e^{-1-v}),\qquad
 b(v)=-W_{-1}(-e^{-1-v}),
\]
where $W_0$ and $W_{-1}$ are the real inverse branches of $w e^w$.
Writing $\overline F_\Gamma(t)=(1+t)e^{-t}$, the canonical deficit
CDF is zero for $v<0$ and is
\begin{equation}
 F_I(v)=\overline F_\Gamma(a(v))-\overline F_\Gamma(b(v)),
                \qquad v\geq0.\label{final:eq-deficit-cdf}
\end{equation}
It has no atom at zero and, for $v>0$, has density
\begin{equation}
 f_I(v)=e^{-1-v}
 \left(\frac{a(v)}{1-a(v)}+\frac{b(v)}{b(v)-1}\right),
 \qquad f_I(v)\sim\frac{\sqrt2}{e\sqrt v}\quad(v\downarrow0).
 \label{final:eq-deficit-density}
\end{equation}
The canonical coordinate involution is
\begin{equation}
 j_I(t)=
 \begin{cases}
 -W_{-1}(-te^{-t}),&0<t<1,\\
 1,&t=1,\\
 -W_0(-te^{-t}),&t>1.
 \end{cases}\label{final:eq-deficit-involution}
\end{equation}
It is a smooth strictly decreasing involution, with $j_I'(1)=-1$,
unique fixed point one, and $I\circ j_I=I$, $p\circ j_I=p$.
\end{proposition}
\begin{proof}
The derivative $I'(t)=(t-1)/t$ and the infinite endpoint limits
give precisely two roots at every positive level.  Rearranging
$I(t)=v$ gives $(-t)e^{-t}=-e^{-1-v}$ and proves the Lambert
formulas.  The sublevel set is $[a(v),b(v)]$, proving the CDF by
integration.  Implicit differentiation gives
$a'=a/(a-1)$ and $b'=b/(b-1)$, while
$p(a)=p(b)=e^{-1-v}$; differentiating the CDF gives its density.
The Taylor expansion $I(1+h)=h^2/2+O(h^3)$ gives
$a(v)=1-\sqrt{2v}+O(v)$ and
$b(v)=1+\sqrt{2v}+O(v)$, yielding the stated asymptotic.
The zero level is the singleton $\{1\}$, which has Gamma mass zero.

Equal-level pairing immediately gives the involution identities.
To verify smoothness at the joining point, write
$I(t)=(t-1)^2 r(t)$ near one, where $r$ is smooth, positive, and
$r(1)=1/2$.  The function
\[
 q_I(t)=\operatorname{sgn}(t-1)\sqrt{2I(t)}
\]
therefore extends smoothly with derivative one at one.  Away from
one its derivative is strictly positive by the sign of $I'$.
The endpoint limits make it a smooth increasing bijection onto
$\mathbb R$, and
$j_I=q_I^{-1}\circ(-\operatorname{id})\circ q_I$.
This proves all the remaining assertions, including $j_I'(1)=-1$.
\end{proof}

\begin{proposition}[Deficit transform, cumulants, and determinacy]
\label{final:P4-cumulants}
The maximal real moment-generating domain of $V$ is $s<1$, and
\begin{align}
 M_V(s)&=e^{-s}\frac{\Gamma(2-s)}{(1-s)^{2-s}},\label{final:eq-deficit-mgf}\\
 K_V(s)&=\gamma_Es+
 \sum_{n=2}^\infty\frac{1}{n}
       \left(\zeta(n)-\frac{1}{n-1}\right)s^n\quad(|s|<1).
       \label{final:eq-deficit-cgf}
\end{align}
Consequently
\begin{equation}
 \kappa_1(V)=\gamma_E,\qquad
 \kappa_n(V)=(n-1)!
       \left(\zeta(n)-\frac{1}{n-1}\right)\quad(n\geq2).
       \label{final:eq-deficit-cumulants}
\end{equation}
In particular $\operatorname{Var}V=\pi^2/6-1$,
$\kappa_3(V)=2\zeta(3)-1$, and $\kappa_4(V)=\pi^4/15-2$.
Every nonnegative probability law with all finite moments and these
algebraic cumulants equals the law of $V$; no separate exponential
moment assumption on that candidate is necessary.
For $a>0$ and $s<1$, the conditional transform is
\begin{equation}
 \mathbb E[e^{sI(T)}\mid T\geq a]
 =\frac{e^{-s}\Gamma(2-s,(1-s)a)}
       {(1+a)e^{-a}(1-s)^{2-s}},
 \label{final:eq-deficit-truncation}
\end{equation}
where $\Gamma(q,x)=\int_x^\infty u^{q-1}e^{-u}\,du$.
\end{proposition}
\begin{proof}
The identity $p=e^{-1-I}$ gives
\[
 M_V(s)=e^{-s}\int_0^\infty t^{1-s}e^{-(1-s)t}\,dt.
\]
For $s<1$ the Gamma substitution gives
\eqref{final:eq-deficit-mgf}.  At $s=1$ its integrand is the
positive constant $e^{-1}$, and for $s>1$ the tail diverges.
Euler's Gamma product, expanded absolutely for $|s|<1$, gives
\[
 \log\Gamma(1-s)=\gamma_Es+
 \sum_{k=1}^\infty\{-\log(1-s/k)-s/k\}
 =\gamma_Es+\sum_{n=2}^\infty\zeta(n)s^n/n.
\]
Using $\Gamma(2-s)=(1-s)\Gamma(1-s)$ and the logarithm series
in $\log M_V$ proves \eqref{final:eq-deficit-cgf}, and differentiation
gives the cumulants.  The classical evaluations
$\zeta(2)=\pi^2/6$, $\zeta(4)=\pi^4/90$ give the displayed values.

Algebraic cumulants and moments determine each other recursively:
\[
 m_0=1,\qquad
 m_{n+1}=\sum_{j=0}^n\binom nj\kappa_{j+1}m_{n-j}.
\]
This is coefficient comparison in the formal identity $M'=K'M$.
Thus a candidate has the same moments as $V$.  For any $0<s<1$,
monotone convergence of its nonnegative exponential series gives
$\mathbb E e^{sX}=\sum m_ns^n/n!=M_V(s)<\infty$.
Nonnegativity also gives finiteness for negative $s$; uniqueness
from Proposition~\ref{final:P0-uniqueness} finishes the converse.
Restricting the same integral to $[a,\infty)$ and dividing by
$\mathbb P(T\geq a)=(1+a)e^{-a}$ proves the conditional formula.
\end{proof}

\begin{theorem}[The linked deficit law and involution identify the potential]
\label{final:P4}
Let $J:(0,\infty)\to[0,\infty)$ be continuous, satisfy $J(1)=0$,
be strictly decreasing on $(0,1]$ and strictly increasing on
$[1,\infty)$, and tend to infinity at both endpoints.  Suppose
$q_J(t)=e^{-1-J(t)}$ has integral one.  Let $F_J$ be the law of
$J(T_J)$ under this same density $q_J$, and let $j_J$ exchange its
two equal-level roots in the same $t$ coordinate and fix one.
If
\[
 F_J=F_I\quad\hbox{and}\quad j_J=j_I,
\]
with the target values specified in
\eqref{final:eq-deficit-cdf} and \eqref{final:eq-deficit-involution},
then $J=I$ and $q_J=p$.  Conversely $I$ supplies these data.
The full list \eqref{final:eq-deficit-cumulants} may replace $F_J=F_I$.
More generally the complete pair $(F_J,j_J)$ uniquely determines
every potential in this candidate class.
\end{theorem}
\begin{proof}
Let $a_J(v)\leq1\leq b_J(v)$ be the inverse branches, including
$a_J(0)=b_J(0)=1$, and put $w_J(v)=b_J(v)-a_J(v)$.
The pushforward of Lebesgue measure under $J$ has locally finite
Stieltjes distribution $w_J$: the sublevel interval has exactly
that length.  Weighting by $e^{-1-J}$ therefore gives, as measures,
\[
 dF_J(v)=e^{-1-v}\,dw_J(v),\qquad
 w_J(v)=e\int_{[0,v]}e^s\,dF_J(s).
\]
Both widths vanish at zero, so there is no undetermined integration
constant.  The map $D_J(t)=t-j_J(t)$ is a continuous strictly
increasing bijection from $[1,\infty)$ to $[0,\infty)$:
$j_J$ is decreasing, fixes one, and tends to zero as $t\to\infty$.
Consequently
\[
 b_J(v)=D_J^{-1}(w_J(v)),\qquad a_J(v)=j_J(b_J(v)).
\]
The observed pair recovers both inverse branches and hence $J$
pointwise.  This proof uses Stieltjes measures and does not require
a derivative of $J$ or a density of $F_J$.
The cumulant alternative follows from
Proposition~\ref{final:P4-cumulants}, since $J(T_J)\geq0$.
\end{proof}

\begin{proposition}[Neither linked observation alone suffices]
\label{final:P4-boundaries}
Within the normalized candidate class of Theorem~\ref{final:P4},
there are noncanonical potentials with exactly $F_I$, and others
with exactly $j_I$.  Thus both canonical target observations and
their common density--potential linkage carry identifying information.
The law of $I(T)$ alone does not identify an arbitrary source law
for $T$, even when the potential $I$ itself is fixed.
\end{proposition}
\begin{proof}
Start with the canonical inverse branches $a,b$.  Choose a nonzero
smooth function $h$ supported on a compact level interval inside
$(0,\infty)$ and set $a_\varepsilon=a+\varepsilon h$,
$b_\varepsilon=b+\varepsilon h$.  On that compact interval the
strictly negative derivative of $a$ and strictly positive derivative
of $b$ are bounded away from zero, and the branches are separated
from zero and one.  Small $|\varepsilon|$ therefore preserves the
branch ranges and monotonicities.  Their inverses define an
admissible $J_\varepsilon\neq I$.  The width is unchanged; the
Stieltjes formula in Theorem~\ref{final:P4} preserves both the entire
deficit law and its total probability mass.  Joint uniqueness shows
that its involution differs.

Conversely choose nonzero nonnegative smooth level bumps $f,g$ with
disjoint compact supports in $(0,\infty)$, and set
$h_{\varepsilon,\delta}(v)=v+\varepsilon f(v)+\delta g(v)$.
Small coefficients make $h$ a strictly increasing bijection fixing
zero, with $h'(0)=1$.  The potential $h\circ I$ has exactly $j_I$.
Its mass functional
\[
 N(\varepsilon,\delta)=
 \int_0^\infty e^{-1-h_{\varepsilon,\delta}(I(t))}\,dt
\]
is smooth near zero, since the perturbations have compact level
support, and
$\partial_\delta N(0,0)=-\int p(t)g(I(t))\,dt<0$.
The implicit function theorem gives $\delta(\varepsilon)$ with
$N(\varepsilon,\delta(\varepsilon))=1$.  For nonzero small
$\varepsilon$, disjoint supports ensure that this potential differs
from $I$, and joint uniqueness makes its deficit law different.

Finally draw $V$ with law $F_I$ and set $T_+=b(V)$.  Then
$I(T_+)=V$ exactly, but $T_+\geq1$ almost surely, so its law is not
$\mu_\Gamma$.  This explicitly exhibits the lost allocation between
the two branches if the source law is not linked to $e^{-1-I}dt$.
Merely possessing some linked pair, without prescribing its two
canonical values, plainly holds for all the potentials just built.
\end{proof}

\begin{theorem}[Survival, hazard, logistic equation, and causal Green kernel]
\label{final:P5}
The Gamma survival and hazard are
\begin{equation}
 \overline F_\Gamma(x)=(1+x)e^{-x}
 =\frac{p(1+x)}{p(1)}=e^{-I(1+x)},\qquad
 h_\Gamma(x)=\frac{x}{1+x}\quad(x\geq0).
 \label{final:eq-survival}
\end{equation}
Each of the following identifying clauses recovers $p$:
\begin{enumerate}
\item an absolutely continuous probability on $[0,\infty)$ with
      positive survival at finite arguments, survival value one at
      zero, and actual hazard $x/(1+x)$ almost everywhere;
\item that same survival framework, with its marked hazard
      $v(u)=h(e^u)$ differentiable on $\mathbb R$ and satisfying
      $v'=v(1-v)$, $v(0)=1/2$;
\item a $C^2$ function $f$ on $(0,\infty)$ extending continuously
      to zero, with $f''+2f'+f=0$, $f(0)=0$, and
      $\int_0^\infty f(t)\,dt=1$;
\item a distribution $g$ on $\mathbb R$ supported in $[0,\infty)$
      with $(D+1)^2g=\delta_0$.
\end{enumerate}
In the last clause the result is
$g(t)=\mathbf 1_{[0,\infty)}(t)t e^{-t}$.
If $k(t)=\mathbf 1_{[0,\infty)}(t)e^{-t}$, then $g=k*k$.
Thus $\mu_\Gamma$ is also the law of a sum of two independent
unit-rate exponential variables and the second arrival time of a
unit-rate Poisson process.
\end{theorem}
\begin{proof}
Integration by parts gives the survival formula, and its negative
derivative divided by the survival gives the hazard.  For the first
converse, local absolute continuity and positivity imply
\[
 (\log\overline F)'=-\frac{x}{1+x}\quad\hbox{almost everywhere}.
\]
Integration from zero gives
$\log\overline F(x)=-x+\log(1+x)$; differentiating recovers the
density $p$ almost everywhere.  The function
$v(u)=(1+e^{-u})^{-1}$ solves the second clause.  Uniqueness for the
scalar differential equation, whose right-hand side is locally
Lipschitz, gives this solution on all of $\mathbb R$; the substitution
$x=e^u$ returns the first clause.

For the third clause, $(e^tf)''=0$, so
$f(t)=(A+Bt)e^{-t}$.  The endpoint gives $A=0$ and integration
gives $B=1$.  For the distributional clause, multiplication by
$e^t$ conjugates $(D+1)^2$ to $D^2$, while
$e^t\delta_0=\delta_0$.  The function $t_+=\max(t,0)$ has
distributional second derivative $\delta_0$.  The difference
$e^tg-t_+$ has second derivative zero, hence is an affine
distribution; its support in $[0,\infty)$ forces both affine
coefficients to vanish.  This proves existence and uniqueness.
Direct convolution gives
\[
 (k*k)(t)=\mathbf 1_{[0,\infty)}(t)
       \int_0^t e^{-s}e^{-(t-s)}\,ds
       =\mathbf 1_{[0,\infty)}(t)t e^{-t}.
\]
The convolution is the independent-sum density.  For a unit-rate
Poisson process the event that its second arrival exceeds $x$ is
the event of at most one arrival by $x$, with probability
$e^{-x}(1+x)$.  This gives the final construction.
\end{proof}

\begin{proposition}[Precise survival boundaries and Gamma shift uniqueness]
\label{final:P5-survival-boundaries}
Local absolute continuity and the survival anchor in
Theorem~\ref{final:P5} cannot be dropped from its almost-everywhere
hazard formulation.  For the logarithmic random variable $Y=\log T$,
the hazard with respect to its own coordinate is
\[
 h_Y(u)=e^u h_\Gamma(e^u)=\frac{e^{2u}}{1+e^u},
\]
which differs from the composition $v(u)$ in that theorem.
For a Gamma density of shape $\alpha>0$ and rate $\beta>0$, the
self-survival shift identity
\begin{equation}
 \overline F_{\alpha,\beta}(x)
 =\frac{f_{\alpha,\beta}(x+c)}{f_{\alpha,\beta}(c)}
                  \quad(x\geq0),\qquad c>0,
 \label{final:eq-gamma-self-shift}
\end{equation}
holds exactly when $\alpha=1$ with arbitrary $c$, or when
$\alpha=2$ and $c=1/\beta$.
\end{proposition}
\begin{proof}
The mixtures $r\delta_0+(1-r)\mu_\Gamma$, $0<r<1$, retain the
hazard at positive arguments and change the survival anchor.
For the regularity boundary, let $C$ be a nonconstant continuous
nondecreasing singular function on $[0,\infty)$, zero on $[0,1]$,
constant on $[2,\infty)$, and with derivative zero almost everywhere.
For instance, translate a Cantor distribution function to $[1,2]$.
Then
\[
 \overline F(x)=(1+x)e^{-x}e^{-C(x)}
\]
is a continuous decreasing survival with value one at zero and
limit zero, and satisfies the same logarithmic derivative equation
almost everywhere.  Its singular component makes it different
from $\overline F_\Gamma$.  This example concerns a bare derivative
equation without absolute continuity.  For $Y=\log T$, the density
is $e^u p(e^u)$ and the survival is
$\overline F_\Gamma(e^u)$, which gives the asserted Jacobian factor.

For \eqref{final:eq-gamma-self-shift}, the density is
$f_{\alpha,\beta}(x)=\beta^\alpha x^{\alpha-1}e^{-\beta x}/
\Gamma(\alpha)$.  If $0<\alpha<1$, the right derivative of the
survival at zero is $-\infty$, while the shifted ratio has finite
derivative, so equality is impossible.  If $\alpha=1$, the ratio
and survival are both $e^{-\beta x}$ for every $c$.
For $\alpha>1$, the survival derivative at zero is zero; equality
forces $(\alpha-1)/c-\beta=0$.  For that $c$, the logarithmic
derivative of the ratio is $-\beta x/(x+c)$, so its derivative is
$-(\beta/c)x(1+o(1))$ near zero.  The survival derivative is
$-\beta^\alpha x^{\alpha-1}(1+o(1))/\Gamma(\alpha)$.
Their equality forces $\alpha-1=1$, then $c=1/\beta$.
The Gamma shape-two survival is $(1+\beta x)e^{-\beta x}$,
which verifies the surviving case.  In particular this shift
property without a specified comparison class is not an
unrestricted characterization of $p$.
\end{proof}

\begin{proposition}[Probability Stieltjes transform and its inverse]
\label{final:P5-stieltjes}
For $z>0$,
\begin{equation}
 G_\Gamma(z)=\int_0^\infty\frac{p(t)}{t+z}\,dt
 =1-z e^z E_1(z),\qquad
 E_1(z)=\int_z^\infty\frac{e^{-u}}u\,du.
 \label{final:eq-gamma-stieltjes}
\end{equation}
The complete transform, or its values on any nonempty open interval
in $(0,\infty)$, identifies a candidate finite positive measure on
$[0,\infty)$ as $\mu_\Gamma$.
The Laplace transform $(1+\lambda)^{-2}$ is completely monotone
but is not a nonzero Stieltjes function.  Neither $p$ nor
$\overline F_\Gamma$ is completely monotone.
\end{proposition}
\begin{proof}
Split $t/(t+z)=1-z/(t+z)$ and substitute $u=t+z$ in the second
integral; this proves \eqref{final:eq-gamma-stieltjes}.
For a finite measure $\mu$, Tonelli's theorem gives
\[
 G_\mu(z)=\int_0^\infty e^{-zv}L_\mu(v)\,dv,
 \qquad L_\mu(v)=\int e^{-vt}\,\mu(dt).
\]
The Stieltjes transform is holomorphic for $\operatorname{Re}z>0$,
by local dominated differentiation.  Equality on an open positive
interval thus extends to all positive $z$.  For two such measures
let $h=L_\mu-L_\nu$, a bounded continuous function.  The finite signed
measure $e^{-v}h(v)\,dv$ has zero Laplace transform by the displayed
identity evaluated at $z+1$.  Proposition~\ref{final:P0-uniqueness}
makes it zero; continuity gives $h=0$, and a second application of
Laplace uniqueness gives $\mu=\nu$.

Complete monotonicity of $(1+\lambda)^{-2}$ follows by differentiating
its positive Laplace integral.  A Stieltjes representation is
\[
 f(\lambda)=a/\lambda+b+
 \int_{[0,\infty)}\frac{\rho(dt)}{\lambda+t},
 \qquad a,b\geq0,\quad
 \int\frac{\rho(dt)}{1+t}<\infty.
\]
If this function is nonzero, either $a>0$, $b>0$, or $\rho$ has
positive mass on some bounded interval.  In every case
$\liminf_{\lambda\to\infty}\lambda f(\lambda)>0$.
For $(1+\lambda)^{-2}$ that limit is zero, a contradiction.
Finally $p'(t)=(1-t)e^{-t}>0$ for $0<t<1$, while
$\overline F_\Gamma''(t)=(t-1)e^{-t}<0$ there.  These signs violate
the respective complete-monotonicity inequalities.
\end{proof}

\begin{proposition}[Reverse size bias, equilibrium, and marked tilting]
\label{final:P5-transforms}
Let $\mu$ be a probability on $[0,\infty)$ with finite positive
mean $m$.
\begin{enumerate}
\item Its size-biased probability is $\beta(dt)=t\mu(dt)/m$.
      If $\mu(\{0\})=0$, then
      \[
      c=\int_{(0,\infty)}t^{-1}\,\beta(dt)=1/m,
      \qquad \mu(dt)=\frac{\beta(dt)}{ct}.
      \]
      Conversely any probability $\beta$ on $(0,\infty)$ with
      $0<c<\infty$ in this formula has that unique positive inverse.
      In particular the positive size-bias inverse of
      $\mu_\Gamma$ is $\operatorname{Exp}(1)$.
\item Its equilibrium law has canonical decreasing right-continuous
      density $q_E(t)=\mathbb P(X>t)/m$ for $t\geq0$.
      If $\mathbb P(X>0)=1$, its inverse is
      \[
      m=1/q_E(0+),\qquad\mathbb P(X>t)=m q_E(t).
      \]
      In particular the full equilibrium density
      $(1+t)e^{-t}/2$ identifies $\mu_\Gamma$, without assuming a
      density for the candidate original lifetime.
\item For a probability $\mu$ on $\mathbb R$, a marked real tilt
      $\theta$ with $0<M(\theta)=\int e^{\theta t}\mu(dt)<\infty$
      gives $\mu_\theta(dt)=e^{\theta t}\mu(dt)/M(\theta)$ and has
      the inverse
      \[
      \mu(dt)=\frac{e^{-\theta t}\mu_\theta(dt)}
                    {\int e^{-\theta y}\mu_\theta(dy)}.
      \]
      For $\mu_\Gamma$, the admissible real tilts are $\theta<1$,
      and
      \[
      p_\theta(t)=(1-\theta)^2t e^{-(1-\theta)t},\quad
      \kappa_n(T_\theta)=\frac{2(n-1)!}{(1-\theta)^n}.
      \]
      Its mean is $2/(1-\theta)$ and variance
      $2/(1-\theta)^2$.
\end{enumerate}
Allowing an atom at zero destroys the first two inverses.  Forgetting
the tilt mark leaves a rate ambiguity in the third.
\end{proposition}
\begin{proof}
In the first clause, integrating $1/t$ against the definition gives
$c=\mu((0,\infty))/m=1/m$.  The inverse has total mass one and
mean $1/c$ by direct integration, and size-biasing it returns
$\beta$.  For $\beta(dt)=p(t)dt$, one has $c=1$, and division by
$t$ gives $e^{-t}dt$.

The tail integral identity $\int_0^\infty\mathbb P(X>t)dt=m$
normalizes the equilibrium density.  Positivity of the original
lifetime gives its right limit at zero as $1/m$; the displayed
inverse therefore recovers the entire survival, including all
positive atoms.  Two right-continuous monotone representatives
that agree almost everywhere agree everywhere, since a strict
difference at one point would persist on a short interval to its
right.  Thus equality of equilibrium laws suffices.  For the target
density $q_E(0+)=1/2$, the recovered survival is
$(1+t)e^{-t}$, whose negative derivative is $p(t)$.

Multiplication of the tilt formula by $e^{-\theta t}$ and
normalization proves the third inverse.  The Gamma transform gives
$M(\theta)=(1-\theta)^{-2}$ for $\theta<1$ and divergence otherwise.
Its tilted log MGF is
$-2\log(1-u/(1-\theta))$, proving the cumulants.

Every $r\delta_0+(1-r)\operatorname{Exp}(1)$, $0\leq r<1$, has
the same size bias $\mu_\Gamma$.  Similarly
$r\delta_0+(1-r)\mu_\Gamma$ has the same equilibrium law for all
$r<1$, because its tail and mean are both multiplied by $1-r$.
For the tilt boundary, a Gamma law of any rate $b>0$, tilted by
$\theta=b-r$, has the fixed output rate $r>0$.  Hence the output
without its input tilt mark does not recover $b$.
\end{proof}

\begin{corollary}[Equilibrium fixed points]
\label{final:P5-equilibrium-fixed}
A positive lifetime with finite mean $m$ equals its equilibrium law
if and only if it is exponential with rate $1/m$.
\end{corollary}
\begin{proof}
Equality to the equilibrium law makes the original distribution
absolutely continuous with density $\overline F/m$.
Consequently $\overline F'=-\overline F/m$ almost everywhere and
$\overline F(0)=1$.  Integration gives $\overline F(t)=e^{-t/m}$.
This exponential law satisfies the same equation and mean by direct
integration.
\end{proof}

\begin{theorem}[Rooted coefficients, Borel probabilities, and the inverse germ]
\label{final:P6}
There is a unique formal power series $R(z)\in z\mathbb R[[z]]$
with $R=ze^R$, and for $n\geq k\geq1$,
\begin{equation}
 [z^n]R(z)^k=\frac{k}{n}\frac{n^{n-k}}{(n-k)!}.
 \label{final:eq-rooted-coefficients}
\end{equation}
In particular
\[
 R(z)=\sum_{n\geq1}\frac{n^{n-1}}{n!}z^n,
 \qquad \operatorname{radius}(R)=e^{-1},
\]
and its compositional inverse germ is $z=t e^{-t}=p(t)$.
The marked Borel PGF and its probabilities are
\begin{equation}
 B(s)=R(s/e),\qquad
 b_n=\frac{e^{-n}n^{n-1}}{n!}\quad(n\geq1),\qquad
 B(s)=s e^{B(s)-1},\quad B(1)=1.
 \label{final:eq-borel}
\end{equation}
Its inverse germ is $s=t e^{1-t}=e p(t)$, with the factor $e$
part of the marking.  For $k$ independent Borel variables their
sum has probabilities
\begin{equation}
 [s^n]B(s)^k=
 \frac{k}{n}e^{-n}\frac{n^{n-k}}{(n-k)!}\qquad(n\geq k).
 \label{final:eq-borel-forest}
\end{equation}
The complete rooted EGF or complete Borel PGF thus determines the
inverse analytic germ.  If a candidate global density is real
analytic on the connected interval $(0,\infty)$ and agrees with
that inverse germ near zero, it equals $p$ everywhere.  Alternatively
$t e^{-t}$ specifies a canonical global extension of the germ.
\end{theorem}
\begin{proof}
Successive coefficient comparison in $R=ze^R$ determines each
coefficient from preceding ones and gives a unique formal solution
with linear coefficient one.  The inverse function theorem applied
to $t e^{-t}$ at zero gives a convergent solution germ, so its
coefficients are those of the formal solution.  For completeness,
the residue change of variables $z=u e^{-u}$ gives
\begin{align*}
 [z^n]R(z)^k
 &=\operatorname*{Res}_{u=0}
       u^{k-n-1}e^{nu}(1-u)\,du\\
 &=[u^{n-k}]e^{nu}-[u^{n-k-1}]e^{nu}
 =\frac{k}{n}\frac{n^{n-k}}{(n-k)!}.
\end{align*}
When $n=k$, the second coefficient is zero and the same formula
gives one.  The coefficient ratio for $k=1$ tends to $e$,
since it is $(1+1/n)^{n-1}$, so the radius is $e^{-1}$.
The inverse identity and Borel scaling follow by substitution.

For $0\leq s<1$, the positive-coefficient solution is the inverse
on $[0,1)$ of the strictly increasing map
$t\mapsto t e^{1-t}$ on $[0,1)$.  Indeed the local inverses agree,
and continuation along this interval is possible because its
derivative is $e^{1-t}(1-t)>0$.  Hence $B(s)\uparrow1$ as
$s\uparrow1$.  Monotone convergence of the nonnegative series
proves $\sum b_n=1$.  Products of PGFs for independent sums and
\eqref{final:eq-rooted-coefficients} give
\eqref{final:eq-borel-forest}.

For labelled rooted trees, removing the root leaves an unordered
set of rooted trees on disjoint sets of labels.  The exponential
formula gives precisely $R=ze^R$; hence their number on $n$ labels
is $n^{n-1}$.  An ordered $k$-tuple of such components has EGF
$R^k$; an unordered set of $k$ components has EGF $R^k/k!$.
These conventions are part of the respective count statements.
Finally equality of the inverse germs and the real-analytic
identity theorem prove the global converse in its stated class.
\end{proof}

\begin{theorem}[Borel total size identifies a one-ancestor iid GW law]
\label{final:P6-gw}
In a Galton--Watson tree started from one ancestor, with independent
identically distributed offspring counts and the usual rooted-tree
construction fixed, the total size has law \eqref{final:eq-borel}
if and only if the offspring law is $\operatorname{Pois}(1)$.
This determines the full tree law within that construction.
\end{theorem}
\begin{proof}
Construct a countable independent family of offspring counts indexed
by all finite words in positive integers, and retain a word as a
vertex exactly when its ancestors and the indicated child indices
are present.  For Poisson offspring the offspring PGF is
$\Phi(w)=e^{w-1}$.  The extinction probability is the increasing
limit of its iterates starting at zero, and therefore the least
fixed point of $\Phi$ in $[0,1]$.  For $0<w<1$ one has
$\log w<w-1$, so $\Phi(w)>w$; also $\Phi(0)>0$.
Its only such fixed point is one.  Extinction is therefore almost
sure, and the tree, having finite offspring counts, has finite
total size almost surely.

For any offspring PGF $\Phi$ in this iid construction, conditioning
on the root offspring and using independence of its descendant
trees gives $B(s)=s\Phi(B(s))$ for the total-size PGF when the total
size is almost surely finite.  For the forward Poisson construction
this is $B=se^{B-1}$, whose formal coefficients are the Borel
probabilities in Theorem~\ref{final:P6}.  Conversely, if the observed
total-size law is Borel, then its PGF satisfies both equations.
The function $B$ maps $(0,1)$ onto $(0,1)$, so division by $s$ gives
$\Phi(w)=e^{w-1}$ for every $0<w<1$.  Analytic coefficient uniqueness
identifies every offspring probability as $e^{-1}/n!$.
The independent product construction then fixes the tree law.
\end{proof}

\begin{proposition}[Tree and continuation boundaries]
\label{final:P6-boundaries}
A complete Borel size law does not identify an arbitrary random
rooted tree.  A complete inverse germ does not identify an arbitrary
smooth positive probability density on $(0,\infty)$, even with the
canonical maximum, strict two-branch shape, and endpoint limits.
Criticality or a leading Borel tail asymptotic alone does not
identify the complete Borel law.
\end{proposition}
\begin{proof}
Sample $N$ with probabilities $b_n$ and return either a rooted path
with $N$ vertices or a rooted star with $N$ vertices.  These laws
have the same size, but on $N=3$ their root degrees differ, and
$b_3>0$.  A Poisson GW tree also has positive probability of root
degree two, which a path never has.  Likewise the PGF product in
\eqref{final:eq-borel-forest} constructs independent components;
the sizes of an independently designated forest need not be
independent merely because component marginals have been specified.

Choose nonzero $h\in C_c^\infty((2,3))$ with $\int h=0$.
For small nonzero $\varepsilon$, set $p_\varepsilon=p+\varepsilon h$.
On the compact support both $p$ and $-p'$ have positive minima,
so positivity and strict decrease survive for small $|\varepsilon|$.
Outside the support it equals $p$.  The integral remains one, the
complete germ at zero is unchanged, and the maximum at one and
all endpoint limits remain canonical.  Nevertheless the density
is different globally.

The critical offspring law assigning probabilities $1/2,1/2$ to
zero and two children has mean one and is not Poisson; its trees
have only odd total sizes.  Finally Stirling's formula gives
$b_n\sim(2\pi)^{-1/2}n^{-3/2}$.  Altering $b_1,b_2$ by opposite
sufficiently small amounts preserves positivity, total mass, and
this exact tail asymptotic, while changing the probability law.
\end{proof}

\begin{theorem}[The calibrated maximum-entropy characterization]
\label{final:P7}
Among Lebesgue probability densities $f$ on $(0,\infty)$ satisfying
\[
 \int t f(t)\,dt=2,\qquad
 \int f(t)\log t\,dt=1-\gamma_E,
 \qquad \int f(t)|\log t|\,dt<\infty,
\]
the unique maximizer of differential entropy is $p$, up to equality
almost everywhere, and its entropy is $1+\gamma_E$.
Here the extended convention allows entropy $-\infty$ and defines
it by the relative-entropy formula below; for finite entropy this
agrees with $-\int f\log f$.
\end{theorem}
\begin{proof}
The target constraints follow from Theorem~\ref{final:P1}.
Since $\log p=\log t-t$, every candidate has the finite cross entropy
$\int f\log p=-1-\gamma_E$.  Put $z=f/p$ and define
\[
 D(f\Vert p)=\int_0^\infty p(t)
          \{z(t)\log z(t)-z(t)+1\}\,dt\in[0,\infty].
\]
The integrand is nonnegative and vanishes exactly when $z=1$.
Because both densities have mass one, this is the usual relative
entropy, and
$\mathsf h(f)=1+\gamma_E-D(f\Vert p)\leq1+\gamma_E$.
For $p$ the divergence is zero; equality for a candidate forces
$f=p$ almost everywhere.  This proves both maximality and its exact
converse.  The two constraints alone do not identify $p$: use four
disjoint compactly supported smooth bumps, solve the three linear
orthogonality equations against $1,t,\log t$, and add a sufficiently
small nonzero resulting perturbation to $p$.  Positivity and all
three equations survive, while the strict entropy inequality holds.
\end{proof}

\begin{theorem}[Unique probability convolution completion]
\label{final:P8}
In the category of probabilities on $[0,\infty)$, there is a unique
family $(\mu_r)_{r\geq0}$ with $\mu_0=\delta_0$,
$\mu_{r+s}=\mu_r*\mu_s$, and $\mu_1=\mu_\Gamma$.
It is
\begin{equation}
 \mu_r=\operatorname{Gamma}(2r,1)\ (r>0),\qquad
 L_{\mu_r}(\lambda)=(1+\lambda)^{-2r}\quad(\lambda\geq0).
 \label{final:eq-gamma-semigroup}
\end{equation}
Weak continuity in $r$ is a conclusion, not an additional premise.
The associated process with stationary independent increments is uniquely
specified in finite-dimensional law once that process category and
the starting value zero are supplied.
\end{theorem}
\begin{proof}
Euler's Gamma integral gives the displayed transforms, and their
products prove the convolution equations by Laplace uniqueness.
For uniqueness let $(\eta_r)$ be another family.  At fixed
$\lambda\geq0$, its Laplace transform is strictly positive and at
most one, so $g(r)=-\log L_{\eta_r}(\lambda)$ is a nonnegative
additive function of $r\geq0$.  It is monotone, since
$g(s)-g(r)=g(s-r)\geq0$ when $s\geq r$.  Rational additivity and
squeezing by rational times yield
$g(r)=r g(1)=2r\log(1+\lambda)$ for every real $r\geq0$.
Another application of Laplace uniqueness gives \eqref{final:eq-gamma-semigroup}.
The continuity theorem for probability Laplace transforms now
gives weak continuity, since these transforms vary continuously
and their pointwise limits are the displayed probability transforms.

For increasing times, take independent increments with the laws
indexed by the successive time differences and form their partial
sums.  Convolution ensures consistency under inserting or deleting
times; the probability extension theorem constructs the process
law, and these prescribed products determine every finite-dimensional
distribution.  In the usual subordinator path convention one takes
its right-continuous increasing version.  The time-one marginal
alone does not specify an independently supplied process: the
process $X_r=rT$ with $T\sim\mu_\Gamma$ has that marginal and
dependent increments, since two unit increments are the same
nonconstant variable $T$.
\end{proof}

\begin{proposition}[Gamma L\'evy measure, drift, and complete Bernstein structure]
\label{final:P8-levy}
The Gamma convolution semigroup has exponent and L\'evy measure
\begin{equation}
 \psi(\lambda)=2\log(1+\lambda)
 =\int_0^\infty(1-e^{-\lambda x})\frac{2e^{-x}}x\,dx,
 \qquad \nu(dx)=\frac{2e^{-x}}x\,dx.
 \label{final:eq-gamma-levy}
\end{equation}
Among probability subordinators with this exact L\'evy measure, the
following are equivalent and each identifies that semigroup:
zero drift; $\psi(\lambda)/\lambda\to0$; time-one support infimum
zero; and time-one mean two.  The L\'evy measure alone leaves arbitrary
nonnegative drift.  Moreover
\begin{equation}
 \frac{2e^{-x}}x=2\int_1^\infty e^{-sx}\,ds,
 \qquad
 \log(1+\lambda)=\int_1^\infty
                 \frac{\lambda}{s(s+\lambda)}\,ds.
 \label{final:eq-complete-bernstein}
\end{equation}
Thus $\psi$ is a complete Bernstein function and
$\psi'(\lambda)=2/(1+\lambda)$ is Stieltjes.
\end{proposition}
\begin{proof}
The integrability condition
$\int(1\wedge x)\nu(dx)<\infty$ follows from boundedness of
$2e^{-x}$ near zero and its exponential tail after division by $x$.
Differentiating the integral in \eqref{final:eq-gamma-levy} with
respect to $\lambda$ gives $2/(1+\lambda)$ by dominated convergence
on compact nonnegative $\lambda$ intervals.  Its value at zero is
zero, so integration proves the exponent formula.

 The L\'evy--Khintchine representation for a subordinator is given in
 \cite[Theorem~3.2, p.~22]{SSV}.  With this measure it is
 $k+d\lambda+2\log(1+\lambda)$, with killing $k\geq0$
and drift $d\geq0$.  Probability mass one excludes killing and
leaves $d\lambda+2\log(1+\lambda)$.  Its time-one law is
$d+T$, by multiplication of Laplace transforms and uniqueness;
its support infimum is $d$, its mean is $d+2$, and the asymptotic
ratio of the exponent to $\lambda$ is $d$.  This proves all the
equivalences and the drift counterexample.  In particular the
complete time-one law in Theorem~\ref{final:P8} already forces
zero drift.

The first identity in \eqref{final:eq-complete-bernstein} is the
elementary exponential integral; the second follows by integrating
$1/s-1/(s+\lambda)$ from one to infinity.  Differentiation of the
first on compact positive $x$ intervals shows that the L\'evy density
is completely monotone.  By definition this makes its Bernstein
exponent complete.  The derivative is the Stieltjes transform of
$2\delta_1$, proving the last assertion.
\end{proof}

\begin{corollary}[Complete integer tails identify the Gamma exponent]
\label{final:P8-samples}
Let $f$ be a Bernstein function.  If, for some integer $N\geq1$,
\[
 f(n)=2\log(1+n)\qquad\hbox{for every integer }n\geq N,
\]
then $f(\lambda)=2\log(1+\lambda)$ for every $\lambda>0$.
\end{corollary}
\begin{proof}
The target is Bernstein by Proposition~\ref{final:P8-levy}.
Apply the complete integer-tail uniqueness theorem
\ref{final:O6} to $f$ and this target.  That theorem derives its
finite compact moment measure from the L\'evy--Khintchine
representation and does not assume a separately given value at zero.
\end{proof}

\begin{proposition}[Explicit Gamma self-decomposition]
\label{final:P8-selfdecomposition}
For $0<c<1$ there exists a nonnegative compound-Poisson residual
$Y_c$, independent of a Gamma copy $T'$, such that
$T\stackrel{d}{=}cT'+Y_c$.  Its L\'evy measure, total jump intensity,
Laplace transform, and mass at zero are respectively
\begin{align}
 \nu_c(dx)&=\frac{2(e^{-x}-e^{-x/c})}{x}\,dx,
 &\nu_c((0,\infty))&=-2\log c,\label{final:eq-residual-levy}\\
 R_c(\lambda)&=\left(\frac{1+c\lambda}{1+\lambda}\right)^2,
 &\mathbb P(Y_c=0)&=c^2.\label{final:eq-residual-transform}
\end{align}
In particular the Gamma law is infinitely divisible and
self-decomposable.  Those two general properties alone do not
identify its shape or rate.
\end{proposition}
\begin{proof}
The density in \eqref{final:eq-residual-levy} is nonnegative.
Since
$e^{-x}-e^{-x/c}=\int_1^{1/c}x e^{-sx}\,ds$, Tonelli gives its
total mass as $2\int_1^{1/c}ds/s=-2\log c$.
A Poisson number of independent positive jumps with normalized law
$\nu_c/(-2\log c)$ therefore constructs $Y_c$.  Its log Laplace
transform is $-\int(1-e^{-\lambda x})\nu_c(dx)$.
Subtracting the two integrals in \eqref{final:eq-gamma-levy}, with
the change of variables for the second, gives
\[
 \int(1-e^{-\lambda x})\nu_c(dx)
 =2\log(1+\lambda)-2\log(1+c\lambda).
\]
Exponentiation proves \eqref{final:eq-residual-transform}; the
mass at zero is the no-jump probability $e^{2\log c}=c^2$.
On a product space take an independent $T'$ with law $\mu_\Gamma$.
The product of the Laplace transforms of $cT'$ and $Y_c$ is
$(1+\lambda)^{-2}$, proving the decomposition by uniqueness.
Infinite divisibility also follows from the $n$-fold convolution
of the time-$1/n$ law in Theorem~\ref{final:P8}.
The identical construction with density
$\alpha e^{-\beta x}/x$, for arbitrary $\alpha,\beta>0$,
proves both properties for every Gamma shape and rate, providing
distinct laws with those same broad properties.
\end{proof}

\begin{theorem}[One linked residual determines a law on the entire real line]
\label{final:P9}
Fix $0\leq c<1$.  Let $Y_c$ have the probability law whose Laplace
transform is \eqref{final:eq-residual-transform}, equivalently the
law of a sum of two independent variables with law
$c\delta_0+(1-c)\operatorname{Exp}(1)$.
For an arbitrary Borel probability $\mu$ on $\mathbb R$, the
following are equivalent:
\begin{enumerate}
\item $\mu=\mu_\Gamma$;
\item there are independent variables $X'$ and $Y_c$, with
      $X'\sim\mu$, such that $cX'+Y_c\sim\mu$;
\item the characteristic function of $\mu$ satisfies
      \[
      \varphi(\xi)=\varphi(c\xi)
          \left(\frac{1-ic\xi}{1-i\xi}\right)^2
                         \qquad(\xi\in\mathbb R).
      \]
\end{enumerate}
No support, density, or moment assumption on $\mu$ is needed.
The residual itself determines $c$ within this family through
$\mathbb P(Y_c=0)=c^2$, so its scale may be recovered from the
complete residual observation.  The source-law linkage is essential;
the contraction endpoint $c=1$ does not identify $\mu$.
\end{theorem}
\begin{proof}
A variable with law $c\delta_0+(1-c)\operatorname{Exp}(1)$ has
Laplace transform
$c+(1-c)/(1+\lambda)=(1+c\lambda)/(1+\lambda)$.
The sum construction proves existence of the residual and its
transform, including $c=0$.  Both summands are nonnegative and
have an atom of mass $c$ at zero, so the sum has an atom of mass $c^2$ at zero.
Its characteristic function is
$((1-ic\xi)/(1-i\xi))^2$, by the same calculation.

Independence gives the factorization in the third clause from the
second.  Conversely that factorization and a product-space
construction, followed by characteristic-function uniqueness, give
the second clause.  Iterating the factorization for any integer
$n\geq1$ telescopes:
\[
 \varphi(\xi)=\varphi(c^n\xi)
       \left(\frac{1-ic^n\xi}{1-i\xi}\right)^2.
\]
The denominators are nonzero for real $\xi$.  Since $c^n\xi\to0$
and every probability characteristic function is continuous at
zero with value one, letting $n\to\infty$ gives
$\varphi(\xi)=(1-i\xi)^{-2}$.  Theorem~\ref{final:P1} and
characteristic-function uniqueness now identify $\mu$ on all of
$\mathbb R$.  For the converse, the Gamma characteristic function
satisfies the factorization by direct cancellation, and independent
variables with the required laws can be constructed on a product
space.  This also covers $c=0$, where $Y_0$ is already Gamma.

Forgetting linkage permits the very same observed $Y_c$ alongside
an unrelated designated $X$ with either Gamma shape two or Gamma
shape three, so it identifies neither such designation.  At $c=1$
the residual is $\delta_0$ and the factorization reduces to
$\varphi(\xi)=\varphi(\xi)$ for every law.  Thus contraction and
the equation linking the original law to this precise residual
are substantive identifying premises.  Independence may be
replaced by the exact factorization as stated, but that
factorization does not follow for an arbitrary supplied coupling.
\end{proof}

\begin{remark}[Formalization]
\label{final:P-formalization-scope}
The probability development, including its measure-theoretic and
analytic statements and counterexamples, is part of the complete
Lean~4 formalization with no admitted proofs. The Bernstein
integer-tail inverse used above is Theorem~\ref{final:O6}.
\end{remark}
